\documentclass[11pt,a4paper]{article}

\usepackage[polish]{babel}
\usepackage[utf8]{inputenc}
\usepackage{polski}
\usepackage[T1]{fontenc}
\usepackage{indentfirst}
\usepackage{wrapfig}    % for wrapping figures, tables
\usepackage{isotope}
\usepackage{hyperref}
\usepackage{amsmath}
\usepackage{amssymb}
\usepackage{bm}
\usepackage{gensymb}
\usepackage{epsfig}
\usepackage{graphics}
\usepackage[shortlabels]{enumitem}
\usepackage{xspace}
\xspaceaddexceptions{[]\{\}}

\def\showsolutions{}  %← show solutions. Comment out to hide solutions

%
%
%fixpagesize
\pagestyle{empty}
\addtolength{\textwidth}{4cm}
\addtolength{\textheight}{6cm}
\addtolength{\evensidemargin}{-3cm}
\addtolength{\oddsidemargin}{-2cm}
\addtolength{\topmargin}{-3cm}
\parindent=0cm

%
%
% Changes figure placing algorithm
\renewcommand{\topfraction}{1}       % maximal fraction of a page allowed for figures
\renewcommand{\textfraction}{0.15}   % minimal number of text for figure-text shared pages
\renewcommand{\floatpagefraction}{0.95} % if two above does not help, this could do the job 
                                        % must be: floatpagefraction < topfraction !!!!
\renewcommand{\textfraction}{0} % minimum fraction of page, which must be  devoted to text
\renewcommand{\topfraction}{1}  % maximum fraction at top, which can be occupied whit floats
\setcounter{totalnumber}{400}   % increase the number of floats for one page
\setcounter{topnumber}{200}     % at all/top/bottom.
\setcounter{bottomnumber}{200}  %

%
%
%small distance in list/item/enum for enumitem package
\setlist[itemize,enumerate]{topsep=0em}
\setlist{noitemsep}

%Nuclear notations: usage \nucl{235}{92}{U}. Math mode optional
\newcommand{\nucl}[3]{\ensuremath{
  \phantom{\ensuremath{^{#1}_{#2}}}
  \llap{\ensuremath{^{#1}}}
  \llap{\ensuremath{_{\rule{0pt}{.75em}#2}}}
  \mbox{#3} } 
} 

\newcounter{zadanie}\newcommand{\zadanie}[1][]{\addtocounter{zadanie}{1} ~\\  {\bf \emph{Exercise \arabic{zadanie} #1 }} \\}
\newcounter{zaddom}\newcommand{\zaddom}[1][]{\addtocounter{zaddom}{1} ~\\  {\bf \emph{Homework \arabic{zaddom} #1 }} \\}

%dynamic solutions hiding
\usepackage{comment}

\newif\ifshowsolutions
\ifdefined\showsolutions
  \showsolutionstrue
\else
  \showsolutionsfalse
\fi

\ifshowsolutions
  \newenvironment{solution}{\vspace{1cm} \par\textbf{Solution.} \vspace{1cm}}{}
\else
  \excludecomment{solution}
\fi
%%%%%%%%%%%%%%%%%%%%%%%%%%%%%%%%%%%%%%%%%%
\begin{document}         

%%%%%%%%%%%%%%%%%%%%%%%%%%%%%%%%%%%%%%%%%%%%%%%%%%%%%
\begin{centering}
  {\bf {\Large
 Introduction to Experimental Particle Physics I \\
  Set II: Reminder of the basics
  }} \\  
\end{centering}

\vspace{1cm}
%%%%%%%%%%%%%%%%%%%%%%%%%%%%%%%%%%%%%%%%%%%%%%%%%%%%%%%
\zadanie[Fermions]

Please recall mathematics relevant for spin 1/2 particles:

\begin{itemize}
    \item write down general form of a spin component (spinor) of the fermion wave function
    \item write down spin projection operators: 
    $\hat{\sigma}_x$, $\hat{\sigma}_y$, $\hat{\sigma}_z$
    \item what is probability to observe positive/negative z component of spin 
    for a spin 1/2 particle in a eigenstate state od $\hat{\sigma}_x$ 
    operator with positive eigenvalue?
    \item write down eigenvalues and eigenvectors of the spin projection operators
    \item write down rotation matrix for rotation around the $x$ axis by an 
          angle $\phi$ 
    \item what is probability to observe positive z component of spin
          of a particle in a state obtained by rotation of the spinor with 
          positive eigenvalue of $\hat{\sigma}_x$ operator by an angle $\phi$ 
          around the $x$ axis? 
\end{itemize}

\begin{solution}

Spin 1/2 particle has two degrees of freedom, which can be represented 
by a two-component spinor:
\begin{align*}
  \Psi = \begin{pmatrix} a \\ b \end{pmatrix}
\end{align*}

where upper component is the amplitude for spin-up state and lower component 
is the amplitude for spin-down state along the z-axis. 
The spinor is normalized such that $|a|^2 + |b|^2 = 1$.

\vspace{0.2cm}

Spin projection operators are given by the Pauli matrices multiplied by $\hbar/2$:
\begin{align*}
  \hat{\sigma}_x = \frac{\hbar}{2}\begin{pmatrix} 0 & 1 \\ 1 & 0 \end{pmatrix}, \quad
  \hat{\sigma}_y = \frac{\hbar}{2}\begin{pmatrix} 0 & -i \\ i & 0 \end{pmatrix}, \quad
  \hat{\sigma}_z = \frac{\hbar}{2}\begin{pmatrix} 1 & 0 \\ 0 & -1 \end{pmatrix}
\end{align*}

\vspace{0.2cm}

Eigenvalues and eigenvectors of the spin projection operators are:
\begin{align*}
  \hat{\sigma}_x: & \quad \lambda = \pm \frac{\hbar}{2}, \quad \text{eigenvectors: } \frac{1}{
\sqrt{2}}\begin{pmatrix} 1 \\ 1 \end{pmatrix}, \frac{1}{\sqrt{2}}\begin{pmatrix} 1 \\ -1 \end{pmatrix} \\
\hat{\sigma}_y: & \quad \lambda = \pm \frac{\hbar}{2}, \quad \text{eigenvectors: } \frac{1}{
\sqrt{2}}\begin{pmatrix} 1 \\ i \end{pmatrix}, \frac{1}{\sqrt{2}}\begin{pmatrix} 1 \\ -i \end{pmatrix} \\
\hat{\sigma}_z: & \quad \lambda = \pm \frac{\hbar}{2}, \quad \text{eigenvectors: } \begin{pmatrix}
  1 \\ 0 \end{pmatrix}, \begin{pmatrix} 0 \\ 1 \end{pmatrix}
\end{align*}

Eigenstate of $\hat{\sigma}_x$ with positive eigenvalue is 
$\frac{1}{\sqrt{2}}\begin{pmatrix} 1 \\ 1 \end{pmatrix}$, which corresponds 
to a superposition of spin-up and spin-down states along the z-axis.
The probability to observe positive z component of spin is:
\begin{align*}
  P(\text{spin up}) = \left| \langle \uparrow_z | \uparrow_x \rangle \right|^2 = \left| \frac{1}{\sqrt{2}} \right|^2 = \frac{1}{2}
\end{align*}
The probability to observe negative z component of spin is:
\begin{align*}
  P(\text{spin down}) = \left| \langle \downarrow_z | \uparrow_x \rangle \right|^2 = \left| \frac{
1}{\sqrt{2}} \right|^2 = \frac{1}{2}
\end{align*}

\vspace{0.2cm}

The rotation matrix for rotation around the $x$ axis by an angle $\phi$ is given by:
\begin{align*}
  R_x(\phi) = \begin{pmatrix} \cos(\phi/2) & -i\sin(\phi/2) \\ -i\sin(\phi/2) & \cos(\phi/2) \end{pmatrix}
\end{align*}

\vspace{0.2cm}

Probability to observe positive z component of spin of a particle in a state 
obtained by rotation of the spinor with positive eigenvalue of $\hat{\sigma}_x$ 
operator by an angle $\phi$ around the $x$ axis is:
\begin{align*}
  P(\text{spin up}) &= \left| \langle \uparrow_z | R_x(\phi) | \uparrow_x \rangle \right|^2 \\
  &= \left| \begin{pmatrix} 1 & 0 \end{pmatrix}
     \frac{1}{\sqrt{2}}
     \begin{pmatrix}
       \cos(\phi/2) \\
       -i\sin(\phi/2)
     \end{pmatrix}
     \right|^2 \\
  &= \frac{1}{2} \left| \cos(\phi/2) \right|^2 
\end{align*}

\end{solution}
%%%%%%%%%%%%%%%%%%%%%%%%%%%%%%%%%%%%%%%%%%%%%%%%%%%%%%%
%%%%%%%%%%%%%%%%%%%%%%%%%%%%%%%%%%%%%%%%%%%%%%%%%%%%%%%
\zadanie[Heavy meson naive model]

Consider a $\Upsilon$ meson consisting of a bottom quark and antiquark.
Assume the quark-quark strong interaction can be (over simplified) described 
by $1/r$ term only: $V(r) = -\frac{4}{3}\alpha_s/r$, where $\alpha_s$ is the strong 
coupling constant. The meson mass can be estimated as a sum of the quark masses and 
the energy of the ground state of the quark-antiquark system in the potential $V(r)$: 

\begin{align*}
M_{\Upsilon} = 2 m_b + E_{1S}
\end{align*}

Take $\alpha_s = 0.2$, $m_b = 4.8$ GeV, and estimate the mass of the $\Upsilon$ meson using 
the Bohr model of the hydrogen atom. Compare your result with the measured mass 
of the $\Upsilon$: $M_{\Upsilon}^{exp} = 9.4603$ GeV.

\begin{solution}

Formula for the energy levels of the of the hydrogen atom is given by:
\begin{align*}
E_n = -\frac{1}{2} \frac{m_e e^4}{(4\pi \epsilon_0)^2 \hbar^2} \frac{1}{n^2} =
-\frac{\alpha^2 m_e c^2}{2 n^2}
\end{align*}

For the $\Upsilon$ meson we have to make following replacements:
\begin{itemize}
    \item $m_e \rightarrow \mu = \frac{m_b}{2}$
    \item $\alpha \rightarrow \frac{4}{3} \alpha_s$
\end{itemize}

and, for the n=1 state, we get:
\begin{align*}
E_{1S} = -\frac{1}{2} \left( \frac{4}{3} \alpha_s \right)^2 \frac{m_b}{2}
= -\frac{8}{9} \alpha_s^2 m_b = -\frac{4}{9} \cdot 0.2^2 \cdot 4.8~\text{GeV} = -0.170~\text{GeV}
\end{align*}

which gives the following estimate for the $\Upsilon$ meson mass:
\begin{align*}
M_{\Upsilon} = 2 m_b + E_{1S} = 2 \cdot 4.8~\text{GeV} - 0.170~\text{GeV} = 9.429~\text{GeV}    
\end{align*}

Unfortunately, this does not work for other states, as the potential is not 
purely Coulomb like, but has a linear term as well.

\end{solution}
%%%%%%%%%%%%%%%%%%%%%%%%%%%%%%%%%%%%%%%%%%%%%%%%%%%%%%%
%%%%%%%%%%%%%%%%%%%%%%%%%%%%%%%%%%%%%%%%%%%%%%%%%%%%%%%



\end{document}
